% Auto-generated by build_latex.py — do not edit by hand.
% Test 1–tagged questions only

\subsection*{\texttt{test1\_sample\_questions\_q1}}
\textit{test1\_sample\_questions.pdf \#1 · test1}\par\smallskip
\begin{quote}
\small\sloppy
Short Answer Questions.\\
\\
    (a) The time between user requests is exponentially distributed with mean 1 minute.\\
        One minute has passed since the last request. What is the expected time to the next\\
        request?\\
    (b) Calculate the variance of the following random variable: P (X = 0) = 1/3, P (X =\\
        1) = 1/3 and P (X = 2) = 1/3.\\
\end{quote}

\subsection*{\texttt{test1\_sample\_questions\_q2}}
\textit{test1\_sample\_questions.pdf \#2 · test1}\par\smallskip
\begin{quote}
\small\sloppy
The time (measured in years) required to complete a software project has a density of the\\
    form                              (\\
                                         kx(1 − x), 0 ≤ x < 1\\
                              f (x) =\\
                                         0,          otherwise\\
   What is the probability that a project is finished in less than four months?\\
\end{quote}

\subsection*{\texttt{test1\_sample\_questions\_q3}}
\textit{test1\_sample\_questions.pdf \#3 · test1}\par\smallskip
\begin{quote}
\small\sloppy
(a) The lifetime of a component is exponentially distributed with mean 1000 hours. What\\
    is the probability a particular component has a lifetime that does not exceed 2000 hours?\\
    (b) Repeat part (a) if the lifetime is uniformly distributed between 500 and 5000 hours.\\
    (c) A random variable X has density f (x) = 3x2 for x between 0 and 1, and 0 otherwise.\\
    Calculate the variance of X.\\
\end{quote}

\subsection*{\texttt{test1\_sample\_questions\_q4}}
\textit{test1\_sample\_questions.pdf \#4 · test1}\par\smallskip
\begin{quote}
\small\sloppy
A discrete random variable takes on the values 1, 2 or 3 with P (X = 1) = 0.5, P (X =\\
    2) = a, P (X = 3) = b. The mean and variance of X are 1.40 and 1.04, respectively. Find\\
    a and b.\\
\end{quote}

\subsection*{\texttt{test1\_sample\_questions\_q5}}
\textit{test1\_sample\_questions.pdf \#5 · test1}\par\smallskip
\begin{quote}
\small\sloppy
(a) A particular disk is busy 30 percent of the time. Each transaction requires 25 disk\\
        accesses and each access takes 25 milliseconds. In transactions per second, what is\\
        the throughput of the disk?\\
    (b) A system has throughput of 3 jobs per minute, and the average number in the system\\
        is seen to be 23. What is the average time in system (from arrival to departure) of\\
        a job?\\
\end{quote}

\subsection*{\texttt{test1\_sample\_questions\_q6}}
\textit{test1\_sample\_questions.pdf \#6 · test1}\par\smallskip
\begin{quote}
\small\sloppy
(a) A system resource executes an average of 10 requests per minute, with an average\\
        service time of 4 seconds per request. Another resource in the same system executes\\
        an average of 5 requests per minute, with an average service time of 11 seconds per\\
        request. Which resource is the bottleneck?\\
    (b) A transaction requires 2 visits on average to Resource A and 1 visit on average to\\
        Resource B. Each visit to Resource A takes 20 seconds, each visit to Resource B\\
        takes 30 seconds. Which Resource is the bottleneck?\\
     (c) I have three resources with utilizations 0.8, 0.3 and 0.9. Given only this information,\\
         which resource would it be most useful to upgrade?\\
\end{quote}

\subsection*{\texttt{test1\_sample\_questions\_q7}}
\textit{test1\_sample\_questions.pdf \#7 · test1}\par\smallskip
\begin{quote}
\small\sloppy
Consider a system with N users, a CPU and two disks. A request from the users first\\
    visits the CPU. After processing at the CPU, with probability 0.8 the request visits disk\\
    A, 0.16 it visits disk B and 0.04 it is completed and returns to the user. After processing\\
    at either disk, the request returns to the CPU. The mean think time at the users is 5\\
    seconds, the mean processing times (per visit) at the CPU, disk A, and disk B are 30, 25\\
    and 40 milliseconds, respectively.\\
    (a) What is the bottleneck device?\\
    (b) What is the maximum possible system throughput (measured at the users)?\\
    (c) Plot an upper bound on the system throughput as a function of N .\\
    (d) What is the maximum possible disk B utilization?\\
\end{quote}

\subsection*{\texttt{test1\_sample\_questions\_q8}}
\textit{test1\_sample\_questions.pdf \#8 · test1}\par\smallskip
\begin{quote}
\small\sloppy
Consider a system with N users, a CPU and two disks. A request from the users first\\
    visits the CPU. After processing at the CPU, with probability 0.48 the request visits disk\\
    A, 0.48 it visits disk B and 0.04 it is completed and returns to the user. After processing\\
    at either disk, the request returns to the CPU. The mean think time at the users is 3\\
    seconds, the mean processing times (per visit) at the CPU, disk A, and disk B are 30, 25\\
    and 40 milliseconds, respectively.\\
    (a) What is the minimum expected response time?\\
    (b) You are asked to give a recommendation for the number of users this system can\\
    support. Provide such a recommendation and justify it.\\
    (c) If you double the processing rate at the bottleneck device, what improvement (in\\
    percentage terms) would you get for the maximum possible system throughput?\\
\end{quote}

\subsection*{\texttt{test1\_sample\_questions\_q9}}
\textit{test1\_sample\_questions.pdf \#9 · test1}\par\smallskip
\begin{quote}
\small\sloppy
A server that is used for compilation jobs is monitored and reveals the following perfor-\\
    mance data (compilations are the only activity performed by the server):\\
\\
    - Average number of active user logins: 230\\
    - Average time to generate a new compilation request: 300 seconds per user\\
    - Average server utilization: 48 percent\\
    - Average CPU processing demand: 0.63 seconds per compilation\\
\\
    (a) What is the system throughput (in compiles per second)?\\
    (b) What is the average compilation time (from compilation submission by a user to\\
    completion)?\\
\end{quote}

\subsection*{\texttt{test1\_sample\_questions\_q10}}
\textit{test1\_sample\_questions.pdf \#10 · test1}\par\smallskip
\begin{quote}
\small\sloppy
Consider a system with three nodes: a user node, a CPU node, and an I/O node. Users\\
    have think times with mean 40 seconds. A request first visits the CPU. After processing\\
    at the CPU, it goes to the I/O with probability 0.9, otherwise the request is completed\\
    and returned to the user. After processing at the I/O, requests are always returned to the\\
    CPU. The total demand per request for the CPU is 4.2 seconds, while the mean processing\\
    times at the I/O are 0.6 seconds per visit.\\
    (a) Determine the bottleneck device and the maximum possible throughput for the system\\
    (measured at the user node).\\
\\
\\
\\
    (b) How many users can this system support?\\
    (c) You have a budget of 20 dollars. You can decrease the mean demand at the CPU by\\
    0.1 seconds times the amount spent in dollars (fractions are allowed). Similarly, you can\\
    decrease the mean time per visit at the I/O by 0.04 seconds times the amount spent in\\
    dollars. If the goal is to maximize the number of users that the system can support, how\\
    would you allocate the 20 dollars?\\
\end{quote}

\subsection*{\texttt{test1\_sample\_questions\_q11}}
\textit{test1\_sample\_questions.pdf \#11 · test1}\par\smallskip
\begin{quote}
\small\sloppy
A database server receives requests from 50 clients. Each request to the DB server requires\\
    that five reads be made on average from the server’s single disk. The average read time\\
    per record is 9 msec. Each database request requires on average 15 msec of CPU time to\\
    be processed.\\
    (a) What is the bottleneck in the DB server?\\
    (b) With the 50 clients, what is the throughput of the DB server?\\
    (c) Consider the following three proposed design changes: (i) the average number of disk\\
    reads per access can be reduced from 5 to 2.5; (ii) the disk is replaced by one 60 percent\\
    faster; (iii) the CPU speed is doubled. Budgetary constraints allow for the implementation\\
    of two of these three. If the goal is to maximize the system capacity (maximum possible\\
    throughput), which two would you choose? Justify your answer.\\
\end{quote}

\subsection*{\texttt{test1\_sample\_questions\_q12}}
\textit{test1\_sample\_questions.pdf \#12 · test1}\par\smallskip
\begin{quote}
\small\sloppy
A computer system has one CPU and two disks. After monitoring the system for two\\
    hours, the following observations were made. The utilization of the CPU was 43 percent\\
    and the utilization of the first disk was 66 percent. Each transaction to the system makes\\
    5 I/O requests to the first disk and 6 to the second disk. The average processing time\\
    (per request) at the two disks are both 20 msec.\\
\\
     (a) What was the throughput of the system (in transactions per second)?\\
     (b) What was the utilization of the second disk?\\
     (c) What is the total processing demand (in milliseconds per transaction) at the CPU?\\
\end{quote}

\subsection*{\texttt{test1\_sample\_questions\_q13}}
\textit{test1\_sample\_questions.pdf \#13 · test1}\par\smallskip
\begin{quote}
\small\sloppy
You have three devices, A, B and C. Device A has average demand DA = 4 and devices\\
    B and C are such that DB + DC = 6.\\
    (a) Give all possible values of DB and DC that yield the highest possible maximum\\
    throughput. As part of your answer, give the highest possible maximum throughput.\\
    (b) Give all possible values of DB and DC that yield the lowest possible maximum through-\\
    put. As part of your answer, give the lowest possible maximum throughput.\\
    (c) For a specific choice of values of DB and DC from (a), compute the maximum number\\
    of users that the system can support if the average user think time is Z = 40.\\
\end{quote}

\subsection*{\texttt{test1\_sample\_questions\_q14}}
\textit{test1\_sample\_questions.pdf \#14 · test1}\par\smallskip
\begin{quote}
\small\sloppy
Consider the following database server model. Each request to the server requires on\\
    average five records to be read from the server’s single disk. The average read time per\\
    record is 9 msec. Each database request requires 15 msec of CPU time.\\
    (a) What is the maximum number of users that this system can support if the average\\
    user think time is 100 seconds?\\
    (b) You are given three upgrade options:\\
\\
\\
\\
    (i) More indexes can be built into the database to reduce the average number of reads\\
    per access from 5 to 2.5.\\
    (ii) The disk can be replaced by a disk 60 percent faster, so the average processing time\\
    would drop to 5.63 msec.\\
    (iii) The CPU speed can be doubled, i.e. the processing time per request is reduced to 7.5\\
    msec.\\
\\
\\
    If the goal is to maximize the maximum possible system throughput, which option would\\
    you choose? (You can only choose one option.)\\
\end{quote}

\subsection*{\texttt{test1\_sample\_questions\_q15}}
\textit{test1\_sample\_questions.pdf \#15 · test1}\par\smallskip
\begin{quote}
\small\sloppy
A server system consists of three resources. Each user request has resource demands as\\
    follows: the processing times at resources A, B and C are 20, 30 and 15 msec, respectively.\\
    The average number of visits to A, B and C are 5, 4 and 10, respectively.\\
    (a) You are given the following three upgrade options:\\
\\
    (i) Reduce the average number of visits to resource C to 7.\\
    (ii) Reduce the average processing time at resource B to 25 msec.\\
    (iii) Reduce the average processing time at resource C to 10 msec.\\
\\
    If the design goal is to maximize the number of users that can be supported, which option\\
    would you choose?\\
    (b) If the average number of jobs at resource A is 2 and the average waiting time (from\\
    arrival to departure at resource A) is 50 msec, what is the throughput of the system (in\\
    user requests completed per second)?\\
\end{quote}

\subsection*{\texttt{test1\_sample\_questions\_q16}}
\textit{test1\_sample\_questions.pdf \#16 · test1}\par\smallskip
\begin{quote}
\small\sloppy
Consider a system consisting of a single CPU, one disk, and some terminals (clients). The\\
    average think times for each of the terminals is 60 seconds. Suppose that we are presented\\
    with three configurations:\\
                                      System DCP U Ddisk\\
                                         A       4.6      4.0\\
                                         B       5.1      1.9\\
                                         C       3.1      1.9\\
    (a) For System A, the utilization of the CPU is measured to be 0.85. What is the\\
    utilization of the disk?\\
    (b) For System B, what is the maximum possible throughput of the disk?\\
    (c) If there is a requirement that the system be able to support at least 14 users, which\\
    systems satisfy this requirement?\\
\end{quote}

\subsection*{\texttt{test1\_sample\_questions\_q17}}
\textit{test1\_sample\_questions.pdf \#17 · test1}\par\smallskip
\begin{quote}
\small\sloppy
In a timesharing system, accounting log data produced the following profile for user\\
    programs\\
\\
    • each program requires 5 seconds of CPU time, makes 80 I/O requests to disk A and\\
         100 I/O requests to disk B\\
\\
\\
\\
    • average think-time of a user was 18 seconds\\
    • from the device specifications, it was determined that disk A takes 50 milliseconds to\\
         satisfy an I/O request and disk B takes 30 milliseconds per request\\
    • with 17 users, disk A throughput was observed to be 15.70 I/O requests per second\\
    Find the system throughput and the utilizations of the CPU, disk A and disk B.\\
\end{quote}

\subsection*{\texttt{test1\_sample\_questions\_q18}}
\textit{test1\_sample\_questions.pdf \#18 · test1}\par\smallskip
\begin{quote}
\small\sloppy
Jobs arrive to a system at average rate 15 per hour. On observing the system, the average\\
    number occupying it is 10 jobs. What is the mean response time for a job?\\
\end{quote}

\subsection*{\texttt{test1\_sample\_questions\_q19}}
\textit{test1\_sample\_questions.pdf \#19 · test1}\par\smallskip
\begin{quote}
\small\sloppy
Suppose that a chess player has a 90 percent chance of winning any game in a competition.\\
     (a) What is the probability that she wins exactly 3 games, if she plays a total of 5 games?\\
    (b) Suppose that a match ends when one player wins 3 games over another player. What\\
        is the probability that she wins a match?\\
     (c) What is the probability that she wins a match in exactly three games?\\
\end{quote}

\subsection*{\texttt{test1\_sample\_questions\_q20}}
\textit{test1\_sample\_questions.pdf \#20 · test1}\par\smallskip
\begin{quote}
\small\sloppy
Consider a system with three nodes: a user node, a CPU node, and an I/O node. Users\\
    have think times with mean 40 seconds. A request first visits the CPU. After processing\\
    at the CPU, it goes to the I/O with probability 0.9, otherwise the request is completed\\
    and returned to the user. After processing at the I/O, requests are always returned to\\
    the CPU. The mean total demand (over all visits) per request for the CPU is 4.2 seconds,\\
    while the mean processing times at the I/O are 0.6 seconds per visit.\\
     (a) Determine the bottleneck device and the maximum possible throughput for the sys-\\
         tem (measured at the user node).\\
    (b) How many users can this system support?\\
     (c) You have a budget of 20 dollars. You can decrease the mean demand at the CPU by\\
         0.1 seconds times the amount spent in dollars (fractions are allowed). Similarly, you\\
         can decrease the mean time per visit at the I/O by 0.04 seconds times the amount\\
         spent in dollars. If the goal is to maximize the number of users that the system can\\
         support, how would you allocate the 20 dollars?\\
\end{quote}

\subsection*{\texttt{test1\_sample\_questions\_q21}}
\textit{test1\_sample\_questions.pdf \#21 · test1}\par\smallskip
\begin{quote}
\small\sloppy
(a) Database transactions perform an average of 4.5 disk operations on a database server\\
        with a single disk. The database server was monitored during one hour and during\\
        this period, 7,200 transactions were executed. What is the throughput of the disk\\
        (in transactions per second)? If each disk operation takes 20 msec on average, what\\
        is the disk utilization?\\
    (b) What is the average processing time (per operation) for the disk in (a)?\\
     (c) A corporate portal provides Web services to the company’s employees. On average,\\
         500 employees are online requesting Web services from the portal. An analysis of\\
         the portal’s log reveals that 6,480 requests are processed per hour on average. The\\
         average response time was measured as 5 seconds. What is the average think time\\
         of an employee?\\
\end{quote}

\subsection*{\texttt{test1\_sample\_questions\_q22}}
\textit{test1\_sample\_questions.pdf \#22 · test1}\par\smallskip
\begin{quote}
\small\sloppy
A database server receives requests from 50 clients, each client having think time 2 seconds.\\
    The server has one CPU and two disks. Each request requires an average of six accesses\\
    to the first disk and four accesses to the second disk. The average time per access to each\\
    disk is 10 msec and 12 msec, respectively. Each request requires an average of 30 msec of\\
    CPU time.\\
     (a) What is the bottleneck device?\\
     (b) What is the average response time of the server?\\
     (c) You are allowed to make two of the following three changes: (i) the average time per\\
         access for the first disk can be decreased to 7 msec; (ii) the average time per access\\
         for the second disk can be decreased to 11 msec; (iii) the CPU speed can be doubled,\\
         resulting in average CPU time of 15 msec. Which two would you choose, if the goal\\
         is to maximize the maximum possible throughput? Justify your answer.\\
\end{quote}

\subsection*{\texttt{test1\_sample\_questions\_q23}}
\textit{test1\_sample\_questions.pdf \#23 · test1}\par\smallskip
\begin{quote}
\small\sloppy
A random variable X has density\\
                                              (\\
                                                  x−2 1 ≤ x < c\\
                                    f (x) =\\
                                                  0   otherwise\\
    where c > 0. If E[X] = 2, what is c? For your choice of c, calculate V ar(X).\\
\end{quote}

\subsection*{\texttt{test1\_sample\_questions\_q24}}
\textit{test1\_sample\_questions.pdf \#24 · test1}\par\smallskip
\begin{quote}
\small\sloppy
Suppose you have N identical servers. At a given point of time each server has probability\\
    2/3 of working, independent of all of the other servers. What is the minimum value of N\\
    to guarantee that the probability of at least 2 servers working is at least 2/3? Give both\\
    your value of N and the probability that at least 2 servers are working.\\
\end{quote}

\subsection*{\texttt{test1\_sample\_questions\_q25}}
\textit{test1\_sample\_questions.pdf \#25 · test1}\par\smallskip
\begin{quote}
\small\sloppy
A system operates as follows. There are 30 users, each with average think time 20 seconds.\\
    User requests first go to a CPU, with average processing time 0.05 seconds. From the\\
    CPU, with probability 0.05 a request returns to the user, with probability 0.55 it goes to\\
    Disk A, and with probability 0.40 it goes to Disk B. After visiting either disk, requests\\
    always return to the CPU. The average processing times at Disk A and Disk B are 0.08\\
    seconds and 0.04 seconds, respectively.\\
     (a) What is the bottleneck device?\\
     (b) Is an 8 second average response time feasible? If not, what minimum amount of\\
         CPU speedup is required?\\
\end{quote}

\subsection*{\texttt{test1\_sample\_questions\_q26}}
\textit{test1\_sample\_questions.pdf \#26 · test1}\par\smallskip
\begin{quote}
\small\sloppy
Suppose a continuous random variable X has density function f (x) = k/x2 for x > 2 and\\
    0 otherwise. What is V ar(X)? (As part of your solution, you should solve for k.)\\
\end{quote}

\subsection*{\texttt{test1\_sample\_questions\_q27}}
\textit{test1\_sample\_questions.pdf \#27 · test1}\par\smallskip
\begin{quote}
\small\sloppy
A user’s behaviour in interacting with a database system is as follows (we assume that the\\
    user is always connected). There are three actions: query, modify/delete, add. A query\\
    action is followed by another query action with probability 0.8, by a modify/delete action\\
    with probability 0.18 and by an add operation with probability 0.02. A modify/delete\\
    action is followed by a modify/delete action with probability 0.05, otherwise it is followed\\
    by a query action. Finally, an add action is always followed by a query action.\\
\\
\\
\\
     (a) If the first action is a query, what is the probability that the third action is not a\\
         query?\\
    (b) In steady-state, what is the probability that one sees two query actions in a row?\\
\end{quote}

\subsection*{\texttt{test1\_sample\_questions\_q28}}
\textit{test1\_sample\_questions.pdf \#28 · test1}\par\smallskip
\begin{quote}
\small\sloppy
A simple agent in a computer game operates as follows. At each time-step, the agent\\
    stands with probability 0.6, walks with probability 0.3, or runs with probability 0.1. Its\\
    action at each time-step is independent of its action at all other time-steps.\\
\\
     (a) What is the probability that the agent stands, walks and runs, in any order, in 3\\
         consecutive time-steps?\\
    (b) What is the probability over 10 time steps the agent will walk or run at least once?\\
\end{quote}

\subsection*{\texttt{test1\_sample\_questions\_q29}}
\textit{test1\_sample\_questions.pdf \#29 · test1}\par\smallskip
\begin{quote}
\small\sloppy
A web server system produces the following measurements.\\
\\
      • Each web server request requires an average of 1.2 seconds of CPU time and generates\\
        an average of 40 disk accesses to Disk A and 20 disk accesses to Disk B.\\
      • On average, an access to either Disk A or Disk B requires 10 msec.\\
      • Disk A throughput is measured at 30 disk accesses per second.\\
      • The average response time of Disk A is measured to be 25 msec.\\
\\
     (a) What is the utilization of the CPU?\\
    (b) What is the average number of accesses waiting for Disk A?\\
\end{quote}

\subsection*{\texttt{test1\_sample\_questions\_q30}}
\textit{test1\_sample\_questions.pdf \#30 · test1}\par\smallskip
\begin{quote}
\small\sloppy
A random variable X has density\\
                                            (\\
                                                ce−2x 0 < x ≤ 4\\
                                  f (x) =\\
                                                0     otherwise\\
\\
     (a) What is c?\\
    (b) Calculate P (X > 2|X > 1).\\
     (c) Is this distribution memoryless? Justify your answer.\\
\end{quote}

\subsection*{\texttt{test1\_sample\_questions\_q31}}
\textit{test1\_sample\_questions.pdf \#31 · test1}\par\smallskip
\begin{quote}
\small\sloppy
In a file-sharing system, peers are chosen at random, and there is a probability .25 that\\
    a peer has a requested file. The outcome when choosing a peer is independent of the\\
    outcomes of choosing other peers. Peers are chosen until the file is found.\\
\\
     (a) What is the probability that the file is obtained from the third peer selected?\\
    (b) What is the expected number of peers selected before the file is obtained?\\
     (c) Once the file is found on a peer, the time to download it is exponentially distributed\\
         with mean one minute. What is the probability that the download takes longer than\\
         two minutes?\\
\end{quote}

\subsection*{\texttt{test1\_sample\_questions\_q32}}
\textit{test1\_sample\_questions.pdf \#32 · test1 · T/F}\par\smallskip
\begin{quote}
\small\sloppy
True or False.\\
\\
\\
\\
     (a) The uniform distribution has the memoryless property.\\
     (b) If A and B are independent events, P (B|A) = P (B).\\
     (c) The bottleneck in a system is the device with the slowest processing rate.\\
\end{quote}

\subsection*{\texttt{test1\_sample\_questions\_q33}}
\textit{test1\_sample\_questions.pdf \#33 · test1}\par\smallskip
\begin{quote}
\small\sloppy
A continuous random variable has density\\
                                             (\\
                                                 4 3\\
                                                 15 x   1≤x≤c\\
                                   f (x) =\\
                                                 0      otherwise\\
\\
     (a) What is c?\\
\end{quote}

\subsection*{\texttt{test1\_sample\_questions\_q34}}
\textit{test1\_sample\_questions.pdf \#34 · test1 · T/F}\par\smallskip
\begin{quote}
\small\sloppy
True or False.\\
\\
     (a) The variance of X ∼ Exp(1) is less than or equal to the variance of Y ∼ U [0, 2].\\
     (b) If A and B are independent events, P \{B|A\} = P \{B\}.\\
     (c) If X ∼ Exp(10), then P \{X > 15|X > 5\} = P \{X > 15\}.\\
\end{quote}

\subsection*{\texttt{test1\_sample\_questions\_q35}}
\textit{test1\_sample\_questions.pdf \#35 · test1}\par\smallskip
\begin{quote}
\small\sloppy
Suppose that we have the following data on three resources in a system:\\
\\
       • Resource 1 has visit ratio of 5 and average processing time of 10 msec\\
       • Resource 2 has average demand of 35 msec.\\
       • Resource 3 has utilization of 0.5 and throughput of 1/50 msec.\\
\\
    Which resource is the bottleneck? Justify your answer.\\
\end{quote}

\subsection*{\texttt{test1\_sample\_questions\_q36}}
\textit{test1\_sample\_questions.pdf \#36 · test1}\par\smallskip
\begin{quote}
\small\sloppy
Suppose that we have three requests that we must assign to one of three servers. Each\\
    request is assigned such that it is equally likely to go to each of the servers. The three\\
    requests are assigned independently.\\
\\
     (a) What is the probability that the first server receives no requests?\\
     (b) Given that the second server receives no requests, what is the probability that the\\
         first server receives no requests?\\
\end{quote}

\subsection*{\texttt{test1\_sample\_questions\_q37}}
\textit{test1\_sample\_questions.pdf \#37 · test1}\par\smallskip
\begin{quote}
\small\sloppy
Parts coming down an assembly line are inspected in order. If a part is defective, then with\\
    probability .3, the next one is defective. If a part is not defective, then with probability\\
    .1 the next part is defective. Let Xn be the condition (defective/not defective) of the nth\\
    part.\\
\\
      (a) Give the matrix P which describes the Markov chain X.\\
       (b) We observe the 8th part is defective. What is the probability that the 10th part is\\
    also defective?\\
\end{quote}

\subsection*{\texttt{test1\_sample\_questions\_q38}}
\textit{test1\_sample\_questions.pdf \#38 · test1}\par\smallskip
\begin{quote}
\small\sloppy
An analyst is looking at the trends in a certain market and wants to start with a simple\\
    model: he wants to check simply if is up or down on a given week. He finds that if one\\
    week it is up, it is down the next with probability .2 and if one week it is down, the next\\
    it is up with probability .3.\\
    (a) What is the expected number of weeks in the year that the trend is up?\\
    (b) Given that the trend is up one week, what is the probability that is is up two weeks\\
    later?\\
\end{quote}

\subsection*{\texttt{test1\_sample\_questions\_q39}}
\textit{test1\_sample\_questions.pdf \#39 · test1}\par\smallskip
\begin{quote}
\small\sloppy
An electron can be in one of 3 energy states (1,2,3). The state is observed every second\\
    and modelled as a Markov chain where Xn is the energy state at the nth second. If\\
    in state 1, the next second it will be in state 2 with probability .4 and in state 3 with\\
    probability .2. If in state 2, the next second it will be in state 3 with probability .4 and\\
    in state 1 with probability .2. Finally, if in state 3, it is in state 1 the next second with\\
    probability .3 and in state 3 otherwise.\\
    (a) If at time 0 the electron is in state 1, what is the probability that it will be in state 1\\
    at time 1 second and 2 seconds.\\
    (b) Explain in words the interpretation of P 2 (1, 1). Why is this value different from your\\
    answer to (a)?\\
\end{quote}

\subsection*{\texttt{test1\_sample\_questions\_q40}}
\textit{test1\_sample\_questions.pdf \#40 · test1}\par\smallskip
\begin{quote}
\small\sloppy
The number of people in a waiting room is observed every hour and modelled as a Markov\\
    chain. If there are no people in the room, the next hour there is one with probability .5,\\
    two with probability .2 and zero otherwise. If there is one person in the room, the next\\
    hour there are none with probability .5, two with probability .1 and one otherwise. If\\
    there are two people in the room, the next hour there are none with probability .2, one\\
    with probability .6 and two otherwise. In steady-state, calculate\\
    (a) The average number of people waiting.\\
    (b) The variance of the number of people waiting.\\
\end{quote}

\subsection*{\texttt{test1\_sample\_questions\_q41}}
\textit{test1\_sample\_questions.pdf \#41 · test1}\par\smallskip
\begin{quote}
\small\sloppy
The price of a commodity fluctuates between two values, call them high and low. The\\
    price is observed each week. If one week it is high, the next week it is low with probability\\
    .4. If one week it is low, the next week it is high with probability .1. Let \{Xn \} be the\\
    Markov chain which gives the price at week n.\\
    (a) Suppose the initial distribution is Pr\{X0 = high\} = .2, Pr\{X0 = low\} = .8). Calculate\\
    the distribution at the first week (give Pr\{X1 = high\} and Pr\{X1 = low\}).\\
    (b) If the high price is \$110 and the low price is \$90, what is the expected cost at the first\\
    week (using the µ of (a))?\\
\end{quote}

\subsection*{\texttt{test1\_sample\_questions\_q42}}
\textit{test1\_sample\_questions.pdf \#42 · test1}\par\smallskip
\begin{quote}
\small\sloppy
A worker can do three types of jobs (1, 2, 3). If on a given day she is doing job 1, the\\
    next day she is doing job 1 with probability .3. If on a given day she is doing job 2, the\\
    next day she is doing job 2 with probability .6. If on a given day she is doing job 3, the\\
    next day she is doing job 3 with probability .9. If she changes jobs the following day, it\\
    is equally likely that she changes to each of the other two options.\\
    (a) What proportion of days (in steady-state) is she doing each job?\\
\\
\\
\\
    (b) Suppose on day 1 she is doing job 1. What is the probability that the next 2 days she\\
    does both jobs 2 and 3 (in any order)?\\
\end{quote}

\subsection*{\texttt{test1\_sample\_questions\_q43}}
\textit{test1\_sample\_questions.pdf \#43 · test1}\par\smallskip
\begin{quote}
\small\sloppy
A particle can be in one of two states: energized or deenergized. There are two particles\\
    that act independently of each other. At time n, if a particle is energized, it is energized\\
    at time (n + 1) with probability .8. At time n, if a particle is deenergized, it is energized\\
    at time (n + 1) with probability .1.\\
    (a) Let Xn be the number of particles energized at time n. Give its transition matrix P .\\
    (b) Compute Pr\{X52 = 2|X51 = 0, X50 = 0\}.\\
\end{quote}

\subsection*{\texttt{test1\_sample\_questions\_q44}}
\textit{test1\_sample\_questions.pdf \#44 · test1}\par\smallskip
\begin{quote}
\small\sloppy
A machining center has two machines. The number of machines operating on day n is a\\
    Markov chain with state space \{0, 1, 2\} and transition matrix\\
                                                          \\
                                              .1 .6 .3\\
                                        P =  .1 .6 .3 \\
                                                      \\
                                              0 .1 .9\\
\\
    The profit gained per day is \$100 per machine operating.\\
    (a) In steady-state, find the expected profit per day.\\
    (b) For a flat fee of \$10 per day, one can change P (0, 1) = P (1, 1) = 0.5 and P (0, 2) =\\
    P (1, 2) = .4, with the same profit as in (a). Is the change worthwhile?\\
\end{quote}

\subsection*{\texttt{test1\_sample\_questions\_q45}}
\textit{test1\_sample\_questions.pdf \#45 · test1}\par\smallskip
\begin{quote}
\small\sloppy
A user submits requests to a database server. There are three kinds of requests supported:\\
    query, add/modify and delete. If the previous request was a query, then the next request\\
    is an add/modify request with probability 0.2, a delete request with probability 0.1 and\\
    otherwise it is another query request. An add/modify request is followed by another\\
    add/modify request with probability 0.1, otherwise it is followed by a query request. A\\
    delete request is followed by another delete request with probability 0.1, otherwise it is\\
    followed by a query request.\\
\\
     (a) If the first request is a query, what is the probability that the third request is a\\
         query?\\
     (b) In steady-state, what is the proportion of requests that are queries?\\
     (c) A query requires 100 msec of processing, while add/modify and delete requests re-\\
         quire 200 msec of processing. What is the expected processing time per request (in\\
         steady-state)?\\
\end{quote}

\subsection*{\texttt{test1\_sample\_questions\_q46}}
\textit{test1\_sample\_questions.pdf \#46 · test1}\par\smallskip
\begin{quote}
\small\sloppy
Consider a multiprocessor system with two processors and two memory modules. At each\\
    time point (time is discrete), requests are made in the following manner.\\
\\
     (a) If both processors are accessing different memory modules, at the next time point\\
         each processor makes independent requests for the memory modules: with probabil-\\
         ity 0.4, a request is made for memory module 1, otherwise the request is made for\\
         memory module 2.\\
\\
\\
\\
     (b) If both processors are accessing the same module, only one processor makes a new\\
         access request at the next time step (the other maintains its request for the same\\
         memory module). The request probabilities are the same as above.\\
\\
    This can be modelled as a DTMC with three states. Using such a DTMC, calculate the\\
    steady-state probability that the processors are accessing different memory modules.\\
\end{quote}

\subsection*{\texttt{test1\_sample\_questions\_q47}}
\textit{test1\_sample\_questions.pdf \#47 · test1}\par\smallskip
\begin{quote}
\small\sloppy
A computer program’s execution is modelled as a DTMC. The state is the resource being\\
    used at time n. Suppose that there are three resources: CPU (state 1), Disk A (state 2)\\
    and Disk B (state 3). The probability transition matrix is\\
                                                           \\
                                            0.8 0.1 0.1\\
                                      P =  0.1 0.9 0 \\
                                                       \\
                                            0.1 0 0.9\\
\\
     (a) What is the probability that given the program is initially using the CPU, it is using\\
         the CPU two time units later?\\
     (b) What is the steady-state probability that Disk A is being used?\\
\end{quote}

\subsection*{\texttt{test1\_sample\_questions\_q48}}
\textit{test1\_sample\_questions.pdf \#48 · test1}\par\smallskip
\begin{quote}
\small\sloppy
A user’s behaviour in interacting with a database system is as follows (we assume that the\\
    user is always connected). There are three actions: query, modify/delete, add. A query\\
    action is followed by another query action with probability 0.8, by a modify/delete action\\
    with probability 0.18 and by an add operation with probability 0.02. A modify/delete\\
    action is followed by a modify/delete action with probability 0.05, otherwise it is followed\\
    by a query action. Finally, an add action is always followed by a query action.\\
\\
     (a) If the first action is a query, what is the probability that the third action is not a\\
         query?\\
     (b) In steady-state, what is the probability that one sees two query actions in a row?\\
\end{quote}

\subsection*{\texttt{test1\_sample\_questions\_q49}}
\textit{test1\_sample\_questions.pdf \#49 · test1}\par\smallskip
\begin{quote}
\small\sloppy
Consider a DTMC, \{Xn \} with transition matrix\\
                                                            \\
                                           1/2 1/4 1/4\\
                                     P =  0 1/2 1/2 \\
                                                      \\
                                            1   0   0\\
\\
    The states are \{0, 1, 2\}.\\
\\
     (a) Calculate P (X3 = 0|X1 = 0).\\
     (b) Suppose that X0 = 0. As n → ∞, what is the probability that Xn = 2?\\
\end{quote}

\subsection*{\texttt{test1\_sample\_questions\_q50}}
\textit{test1\_sample\_questions.pdf \#50 · test1}\par\smallskip
\begin{quote}
\small\sloppy
Consider a small system of three web pages, labelled A, B and C. Page A is always the\\
    first page visited. After visiting A, pages B and C are equally likely to be visited next.\\
    After page B is visited, page A is always visited next. After page C is visited, pages A\\
    and B are equally likely to be visited next.\\
\\
\\
\\
     (a) What is the probability that A is the fourth page visited?\\
     (b) In steady-state, what is the probability that A is visited?\\
\end{quote}

\subsection*{\texttt{test1\_sample\_questions\_q51}}
\textit{test1\_sample\_questions.pdf \#51 · test1}\par\smallskip
\begin{quote}
\small\sloppy
Consider a system with two components. We observe the state of the system every hour.\\
    A component operating at hour n has probability 0.1 of being failed at hour n + 1. A\\
    component that is failed at hour n has a probability of operating at hour n + 1 of 0.7.\\
    The components operate independently of each other.\\
\\
     (a) Model this system as a DTMC. Define the states and give the transition matrix.\\
     (b) What is the limiting probability that both components are failed at the same hour?\\
\end{quote}

\subsection*{\texttt{test1\_sample\_questions\_q52}}
\textit{test1\_sample\_questions.pdf \#52 · test1}\par\smallskip
\begin{quote}
\small\sloppy
A board is divided into two. There are four markers on the board. Time is discrete. At\\
    each time point, one marker is chosen at random and moved to the opposite side of the\\
    board. What is the limiting probability that the numbers of markers on both sides of the\\
    board are equal?\\
\end{quote}

\subsection*{\texttt{test1\_sample\_questions\_q53}}
\textit{test1\_sample\_questions.pdf \#53 · test1}\par\smallskip
\begin{quote}
\small\sloppy
Suppose that a system has two identical servers. Only one server operates at a time. At\\
    the end of each day, the operating server is inspected. With probability 0.1 it is taken\\
    offline to be repaired and the other server becomes operational (if it is available). Repairs\\
    take exactly two days and both servers can be concurrently repaired.\\
\\
     (a) Model this system as a DTMC and give the P matrix.\\
     (b) Write down the equations that you would use to solve for the limiting distribution\\
         (do not solve them). In terms of your limiting distribution, what is the probability\\
         that you see two consecutive days where a server is taken offline to be repaired?\\
\end{quote}

\subsection*{\texttt{test1\_sample\_questions\_q54}}
\textit{test1\_sample\_questions.pdf \#54 · test1}\par\smallskip
\begin{quote}
\small\sloppy
Consider the following game. A dart will hit the random point Y on the interval (0, 1)\\
    according to the density fY (t) = 2t. You must guess the value of Y (your guess must be\\
    a constant, not random). You will lose \$2 times the magnitude of your error if Y is to the\\
    left of your guess and will lose \$1 times the magnitude of your error if Y is to the right\\
    of your guess. What is the best guess to minimize your expected loss?\\
\end{quote}

\subsection*{\texttt{test1\_sample\_questions\_q55}}
\textit{test1\_sample\_questions.pdf \#55 · test1}\par\smallskip
\begin{quote}
\small\sloppy
Consider a system of three web pages. There is a cache with room for two web pages. The\\
    Least Recently Used (LRU) algorithm is employed to determine the cache contents - the\\
    two most recently used web pages are placed in the cache. Suppose that the probabilities\\
    that the web pages are accessed are 0.8 for the first page, 0.1 for the second page, and\\
    0.1 for the third page.\\
\\
     (a) Determine an appropriate state space and a corresponding P matrix for a DTMC\\
         that can be used to calculate the probability of a cache miss (a request does not find\\
         its web page in the cache).\\
     (b) In terms of the steady-state probabilities, what is the likelihood of a cache miss?\\
         (There is no need to calculate the steady-state probabilities.)\\
\end{quote}

\subsection*{\texttt{test1\_sample\_questions\_q56}}
\textit{test1\_sample\_questions.pdf \#56 · test1}\par\smallskip
\begin{quote}
\small\sloppy
A database server receives requests from 100 clients. The average think time of a client is\\
    10 seconds. The database server has two disks and one CPU. In the current configuration,\\
    the average number of disk accesses per client request is 20 for the first disk, and 30 for\\
    the second disk. For both disks, the average time per access is 10 msec. For each client\\
    request, the total amount of CPU time required has an average of 240 msec.\\
\\
     (a) Is this number of clients well supported by this system? Justify your answer.\\
     (b) You are presented with two design options, both of equal cost. The first would\\
         balance the loads on the disks (keeping the total load over both disks unchanged),\\
         the second would decrease the demand on the CPU to 200 msec. Which option\\
         would you choose? Justify your choice.\\
\end{quote}

\subsection*{\texttt{test1\_sample\_questions\_q57}}
\textit{test1\_sample\_questions.pdf \#57 · test1}\par\smallskip
\begin{quote}
\small\sloppy
You simultaneously request three different files (f1 , f2 , f3 ). The times to download the\\
    files are independent and exponentially distributed with means 1 minute, 1 minute and\\
    2 minutes, respectively. What is the expected time until two files are downloaded (it can\\
    be any two)?\\
\end{quote}

\subsection*{\texttt{test1\_sample\_questions\_q58}}
\textit{test1\_sample\_questions.pdf \#58 · test1}\par\smallskip
\begin{quote}
\small\sloppy
(a) You simultaneously make two requests for the same file from two different servers.\\
        The download times are independent. The download time from the first server is ex-\\
        ponentially distributed with mean 20 seconds and the downlaod time from the second\\
        server is exponentially distributed with mean 30 seconds. What is the probability\\
        that the download from the first server completes first?\\
     (b) Consider the problem in part (a), but now suppose that the download time from the\\
         first server is exactly 20 seconds (the download time from the second server remains\\
         exponentially distributed with mean 30 seconds). What is the probability that the\\
         download from the first server completes first?\\
\end{quote}

\subsection*{\texttt{test1\_sample\_questions\_q59}}
\textit{test1\_sample\_questions.pdf \#59 · test1}\par\smallskip
\begin{quote}
\small\sloppy
Red arrivals occur to a system according to a Poisson process with rate 1 per second.\\
    Green arrivals occur to the same system according to a Poisson process with rate 4 per\\
    second. The arrival processes are independent.\\
\\
     (a) What is the probability that two consecutive arrivals are red?\\
     (b) What is the expected time of the second arrival to the system (colour does not\\
         matter)?\\
     (c) Suppose that we know that exactly one arrival occurred in the first 20 seconds. What\\
         is the probability that it occurred in the first 10 seconds?\\
\end{quote}

\subsection*{\texttt{test1\_sample\_questions\_q60}}
\textit{test1\_sample\_questions.pdf \#60 · test1}\par\smallskip
\begin{quote}
\small\sloppy
A particular computer model fails according to an exponential distribution with rate 1\\
    per year. If you buy two machines of this model, what is the probability that both are\\
    still operating at the end of one year? (Assume that failures are independent events.)\\
\end{quote}

\subsection*{\texttt{test1\_sample\_questions\_q61}}
\textit{test1\_sample\_questions.pdf \#61 · test1}\par\smallskip
\begin{quote}
\small\sloppy
A series of experiments are run as follows. For each experiment, the only possibilities\\
    are success or failure. For a particular experiment, n, a success occurs with probability\\
    0.8 if the previous two experiments (n − 1 and n − 2) were successes. If at least one\\
\\
\\
\\
    of the two previous experiments was a failure, then a success occurs for experiment n\\
    with probability 0.5. Calculate the steady-state probability that a success occurs for an\\
    experiment.\\
\end{quote}

\subsection*{\texttt{test1\_sample\_questions\_q62}}
\textit{test1\_sample\_questions.pdf \#62 · test1}\par\smallskip
\begin{quote}
\small\sloppy
Suppose that in a sample of 10 chips, two chips are defective. You select three chips\\
    (all at once). Let X be the number of defective chips chosen. Give the probability mass\\
    function of X.\\
\end{quote}

\subsection*{\texttt{test1\_sample\_questions\_q63}}
\textit{test1\_sample\_questions.pdf \#63 · test1}\par\smallskip
\begin{quote}
\small\sloppy
Suppose that a random variable X has density\\
                                               (\\
                                                   x   1 ≤ x ≤ b;\\
                                     f (x) =\\
                                                   0   otherwise.\\
\\
    Find E[X]. (Your answer should be a number.)\\
\end{quote}

\subsection*{\texttt{test1\_sample\_questions\_q64}}
\textit{test1\_sample\_questions.pdf \#64 · test1}\par\smallskip
\begin{quote}
\small\sloppy
Consider a system with 50 users, and a server consisting of a CPU and three disks. The\\
    average response time perceived by a user is 10 seconds. The throughput of the system is\\
    1.5 user requests per second. It is also known that each user request requires an average\\
    of 0.52 seconds of CPU time.\\
     (a) What is the average think time of a user?\\
     (b) What is the CPU utilization?\\
     (c) On average, how many requests is the server processing?\\
\end{quote}

\subsection*{\texttt{test1\_sample\_questions\_q65}}
\textit{test1\_sample\_questions.pdf \#65 · test1}\par\smallskip
\begin{quote}
\small\sloppy
A system sends simultaneous requests to to two servers, where the response times have\\
    means of 3 seconds and 2 seconds, respectively. The potential of sending requests to a\\
    third server is being considered, but only if the response from this server is the first to\\
    be received at least 20 percent of the time. If all response times are independent and\\
    exponentially distributed, what is the maximum mean response time that is feasible for\\
    the third server?\\
\end{quote}

\subsection*{\texttt{test1\_sample\_questions\_q66}}
\textit{test1\_sample\_questions.pdf \#66 · test1}\par\smallskip
\begin{quote}
\small\sloppy
Consider the DTMC \{Xn , n = 0, 1, 2, . . .\} with states \{ 0,1,2 \} and transition matrix\\
                                                               \\
                                           1/2 1/4 1/4\\
                                     P =  1/3 0 2/3 \\
                                                      \\
                                           1/2 1/2 0\\
\\
     (a) Suppose that P \{X0 = 0\} = P \{X0 = 1\} = 1/4. Calculate P \{X0 = 2, X1 = 1, X2 =\\
         0\}.\\
     (b) What is the steady-state probability of being in state 0?\\
\end{quote}

\subsection*{\texttt{test1\_sample\_questions\_q67}}
\textit{test1\_sample\_questions.pdf \#67 · test1}\par\smallskip
\begin{quote}
\small\sloppy
(a) A system searches for a file on a number of hosts. Each host searched has a proba-\\
        bility of 0.25 of having the desired file (independently from the other hosts). First,\\
        calculate the probability that at most 2 hosts need to be searched before finding the\\
        file. Now, suppose that you have not found the file after the first four hosts have been\\
        searched. What is the probability that you have to search at most 2 more hosts?\\
\\
\\
\\
     (b) Let a random variable X have density function\\
                                                  (\\
                                                      3x2 0 ≤ x ≤ k\\
                                      fX (x) =\\
                                                      0   otherwise\\
\\
         Calculate k and E[X].\\
\end{quote}

\subsection*{\texttt{test1\_sample\_questions\_q68}}
\textit{test1\_sample\_questions.pdf \#68 · test1}\par\smallskip
\begin{quote}
\small\sloppy
A DTMC with states \{0, 1, 2\} has transition matrix\\
                                                                 \\
                                           1/3              1/3\\
                                     P =  1/2         0\\
                                                                 \\
                                                                  \\
                                                      3/4    0\\
\\
     (a) Fill in the blanks in P (give the complete P in your answer booklet, do not fill them\\
         in on the test sheet).\\
     (b) Calculate P \{X5 = 0, X6 = 0, X7 = 2|X4 = 0, X3 = 1\}.\\
     (c) Calculate the steady-state probability of being in state 1.\\
\end{quote}

\subsection*{\texttt{test1\_sample\_questions\_q69}}
\textit{test1\_sample\_questions.pdf \#69 · test1}\par\smallskip
\begin{quote}
\small\sloppy
Consider the following transition matrix for a DTMC with states \{0, 1, 2\}:\\
                                                                \\
                                             .4 .2 .4\\
                                       P =  .5 .25 .25 \\
                                                       \\
                                             .4 .5 .1\\
\\
     (a) What is P \{X3 = 0|X1 = 2, X0 = 1\}?\\
     (b) What is the steady-state probability of being in state 0?\\
\end{quote}

\subsection*{\texttt{test1\_sample\_questions\_q70}}
\textit{test1\_sample\_questions.pdf \#70 · test1}\par\smallskip
\begin{quote}
\small\sloppy
Suppose that you have an urn where balls are taken out and replaced. The rules for\\
    operation are as follows.\\
\\
     (i) There are always three balls in the urn.\\
     (ii) Each ball is either red or green.\\
    (iii) At each time step, a ball is randomly chosen and replaced by a new ball. With\\
          probability 1/3 the replacement ball is red, otherwise the replacement ball is green.\\
    (iv) We are unsure as to the initial contents (at time step 0) of the urn, so we assume\\
         that all possible initial contents are equally likely.\\
\\
     (a) Calculate the probability that at time step 2 there is at least one red ball in the urn.\\
     (b) Calculate the limiting probability that the three balls are the same colour.\\
\end{quote}

\subsection*{\texttt{test1\_sample\_questions\_q71}}
\textit{test1\_sample\_questions.pdf \#71 · test1}\par\smallskip
\begin{quote}
\small\sloppy
A common model for software reliability is that faults arise according to a Poisson process.\\
    Suppose that we have a product where faults occur according to a Poisson process with\\
    rate 0.2 per day. With probability 0.05, a fault is major, otherwise it is minor.\\
\\
\\
\\
     (a) What is the expected number of major faults in one month? (Assume that one\\
         month is 30 days.)\\
     (b) Suppose that the last fault occurred 3.60 days ago and that it was a minor fault.\\
         What is the expected time until the next minor fault occurs?\\
\end{quote}

\subsection*{\texttt{test1\_sample\_questions\_q72}}
\textit{test1\_sample\_questions.pdf \#72 · test1}\par\smallskip
\begin{quote}
\small\sloppy
(a) A program is executed three times, where the results of the executions are indepen-\\
        dent of each other. The probability that an execution gives a correct result is 0.9.\\
        Let X be the number of correct executions. Give the distribution (c.d.f.) of X and\\
        its expected value.\\
     (b) A continuous random variable X has distribution U [0, 1]. Calculate P \{X > 0.9|X >\\
         0.5\}.\\
\end{quote}

\subsection*{\texttt{test1\_sample\_questions\_q73}}
\textit{test1\_sample\_questions.pdf \#73 · test1}\par\smallskip
\begin{quote}
\small\sloppy
Arrival rates to a server are modelled as follows. There are three rates: high, medium,\\
    and low. The rates change between these three values as follows. If at time unit n it is\\
    high, then at time unit n + 1 it is medium with probability 0.1, low with probability 0.1,\\
    and it remains high otherwise. If at time unit n it is medium, then at time unit n + 1 it\\
    is high with probability 0.2, low with probability 0.2, and it remains medium otherwise.\\
    Finally, if at time unit n it is low, then at time unit n + 1 it is high with probability 0.25,\\
    medium with probability 0.25, and it remains low otherwise.\\
     (a) If the rate is high at time unit n = 0, what is the probability that in the next four\\
         time units, it is high at both n = 2 and n = 4?\\
     (b) Calculate the stationary probability that the request rate is high.\\
\end{quote}

\subsection*{\texttt{test1\_sample\_questions\_q74}}
\textit{test1\_sample\_questions.pdf \#74 · test1}\par\smallskip
\begin{quote}
\small\sloppy
Parts (a) and (b) are unrelated.\\
     (a) A cluster has five servers, two of which have a required piece of data. If the servers\\
         are queried in random order (a server is not queried twice), what is the expected\\
         number of queries before the required piece of data is found?\\
     (b) A chemical company estimates that the probability that one of its batches has a high\\
         impurity level is 0.01. They have implemented a testing procedure that will falsely\\
         give a reading of a high impurity level with probability 0.05 and falsely indicate a\\
         satisfactory impurity level with probability 0.02. If a test results in a satisfactory\\
         impurity level, what is the probability that the impurity level is actually satisfactory?\\
\end{quote}

\subsection*{\texttt{test1\_sample\_questions\_q75}}
\textit{test1\_sample\_questions.pdf \#75 · test1}\par\smallskip
\begin{quote}
\small\sloppy
A load balancer sends five requests to three servers. Each request is sent independently\\
    of the others. With probability 0.2 it is sent to Server 1, with probability 0.4 it is sent\\
    to Server 2, and with probability 0.4 it is sent to Server 3. Calculate the mean and the\\
    variance of the number of jobs sent to Server 1.\\
\end{quote}

\subsection*{\texttt{test1\_sample\_questions\_q76}}
\textit{test1\_sample\_questions.pdf \#76 · test1}\par\smallskip
\begin{quote}
\small\sloppy
A computer random number generator produces values that follow a uniform distribution\\
    between 0 and 1.\\
     (a) If three values are generated, what is the probability that all of the values are less\\
         than 0.4?\\
\\
\\
\\
     (b) If three values are generated, let X be the number of values that are between 0.2\\
         and 0.8. Find the mean and the variance of X.\\
\end{quote}

\subsection*{\texttt{test1\_sample\_questions\_q77}}
\textit{test1\_sample\_questions.pdf \#77 · test1}\par\smallskip
\begin{quote}
\small\sloppy
A professor is setting a test with T/F questions. A question with a T answer is followed\\
    by another question with a T answer with probability 3/4. A question with an F answer\\
    is followed by a question with an F answer with probability 2/3.\\
\\
     (a) Suppose that the first three questions have T answers. What is the probability that\\
         the fifth question has a T answer?\\
     (b) If the exam has 100 questions, provide a good approximation for how many you\\
         would expect to be T answers.\\
\end{quote}

\subsection*{\texttt{test1\_sample\_questions\_q78}}
\textit{test1\_sample\_questions.pdf \#78 · test1}\par\smallskip
\begin{quote}
\small\sloppy
A client/server system is monitored for one hour. During this time, the utilization of\\
    a certain disk is measured to be 50 percent. Each user request makes an average of\\
    2 accesses to this disk, and each disk access has an average processing time of 25 ms.\\
    Suppose that there are 150 clients, each with average think time of 10 seconds. What is\\
    the average response time of the server?\\
\end{quote}

\subsection*{\texttt{test1\_sample\_questions\_q79}}
\textit{test1\_sample\_questions.pdf \#79 · test1}\par\smallskip
\begin{quote}
\small\sloppy
Suppose that we have a system with 30 users that each have average think time of 30\\
    seconds. They submit requests to a server that consists of three devices: CPU, disk\\
    1 and disk 2. The average processing times are E[SCP U ] = 0.05, E[Sdisk1 ] = 0.08,\\
    E[Sdisk2 ] = 0.04 and the visit ratios are E[VCP U ] = 20, E[Vdisk1 ] = 11, and E[Vdisk2 ] = 8.\\
    Is an average response time of 8 seconds feasible? If not, what changes are required?\\
\end{quote}

\subsection*{\texttt{test1\_sample\_questions\_q80}}
\textit{test1\_sample\_questions.pdf \#80 · test1}\par\smallskip
\begin{quote}
\small\sloppy
A DTMC \{Xn \} with states \{0, 1, 2, 3\} has transition matrix\\
                                                                     \\
                                     0 2/3  1/3      0\\
                                   0   x    0    1/2 + x \\
                               P =\\
                                                         \\
                                   0   0 1/3 − y 2/3 − y \\
                                                          \\
\\
                                    1/2 0    0      1/2\\
\\
    Suppose that X0 = 0. Find the expected number of time units until the DTMC reaches\\
    state 3. You should be able to give an explicit number, in other words your final answer\\
    should not depend on x and y.\\
\end{quote}

\subsection*{\texttt{test1\_sample\_questions\_q81}}
\textit{test1\_sample\_questions.pdf \#81 · test1}\par\smallskip
\begin{quote}
\small\sloppy
You interviewed for a software engineer position at a company. The company has two\\
    types of recruiters:\\
\\
       • Type A: gets back to you in a time that is exponentially distributed with mean 10\\
         days.\\
       • Type B: never gets back to you.\\
\\
    Suppose that you are equally likely to be assigned to the two types of recruiters. What\\
    is the probability that your recruiter was Type B given that you have not heard back in\\
    10 days?\\
\end{quote}

\subsection*{\texttt{test1\_sample\_questions\_q82}}
\textit{test1\_sample\_questions.pdf \#82 · test1}\par\smallskip
\begin{quote}
\small\sloppy
A car insurance company puts policy holders in three categories:\\
\\
      • no discount on premiums\\
      • 25\% discount on premiums\\
      • 50\% discount on premiums\\
\\
    New policy holders start with no discount. Their premiums evolve as follows:\\
\\
      • If they made no claims in a given year, they move up one level of discount in the\\
        next year\\
      • If they make at least one claim in a given year, they move down one level of discount\\
        in the next year\\
      • If they have the maximum discount in a given year and make no claims, they remain\\
        at the maximum discount in the next year\\
      • If they have no discount in a given year and make at least one claim, they remain at\\
        the minimum discount in the next year\\
\\
    The probability that a policy holder does not make a claim in a given year is 3/4.\\
\\
     (a) What is the probability that a new policy holder has no discount three years later?\\
    (b) In steady-state, what is the probability that a policy holder has no discount in a\\
        given year?\\
\end{quote}

\subsection*{\texttt{test1\_sample\_questions\_q83}}
\textit{test1\_sample\_questions.pdf \#83 · test1}\par\smallskip
\begin{quote}
\small\sloppy
A system has M users, each with average think time Z = 20 seconds. A user request first\\
    visits a CPU. After visiting the CPU, with probability 0.55 the request goes to Disk 1;\\
    with probability 0.40 the request goes to Disk 2; otherwise the request is complete and\\
    returned to the user. The average processing times are 0.05 seconds at the CPU; 0.08\\
    seconds at Disk 1; 0.04 seconds at Disk 2.\\
\\
     (a) What is the maximum value of M before the system performance begins to degrade?\\
    (b) Is a 10-second average response time feasible when M = 50? If not, what minimum\\
        amount of CPU speedup is required?\\
\end{quote}

\subsection*{\texttt{test1\_sample\_questions\_q84}}
\textit{test1\_sample\_questions.pdf \#84 · test1}\par\smallskip
\begin{quote}
\small\sloppy
A server is either working or down (and being repaired). If it is working today, then\\
    there is a 95\% chance that it will be working tomorrow (otherwise it is down). If it is\\
    down today, there is a 40\% chance that it will be working tomorrow (otherwise it remains\\
    down). If the server is down on a fourth consecutive day, the server is replaced by a new\\
    one, so on the next day it is working.\\
\\
     (a) Model this system as a DTMC and give the corresponding transition matrix.\\
    (b) How would you determine the steady-state probability that the server is working?\\
        Just give the equations that you would solve and a final answer in terms of the\\
        solution to these equations - there is no need to give an explicit solution.\\
\end{quote}

\subsection*{\texttt{test1\_sample\_questions\_q85}}
\textit{test1\_sample\_questions.pdf \#85 · test1}\par\smallskip
\begin{quote}
\small\sloppy
A job can be one of two types. If the job is of Type 1, its processing time is exponentially\\
    distributed with mean 5 minutes. If the job is of Type 2, its processing time is exponen-\\
    tially distributed with mean 60 minutes. A job is twice as likely to be of Type 1 than of\\
    Type 2. If a job has been in process for 4 minutes and is not yet complete, what is the\\
    likelihood that it is a Type 1 job?\\
\end{quote}

